Everything is computed in Gauge Theory Gravity (Definition \claimref{obj.synthesis.gtg-flat-spacetime-gauge-fields}). The Schwarzschild field of Fact \claimref{pairing.imported-schwarzschild-connection-curvature} in the static gauge is, on a constant orthonormal physical frame $\{e_{a}\}$ and with $F = \sqrt f$, $\hb^{-1}(e_{t}) = F\,e_{0}$, $\hb^{-1}(e_{r}) = e_{1}/F$, $\hb^{-1}(e_{\vartheta}) = r\,e_{2}$, $\hb^{-1}(e_{\varphi}) = r\sin\vartheta\,e_{3}$, where $e_{t}, e_{r}, e_{\vartheta}, e_{\varphi}$ are the background frame of the labels; $\Omega(e_{c})$ means $\Omega(\hb(e_{c}))$.

\emph{The gauge field from the torsion equation.} The torsion equation of the definition, read on the constant physical frame, is $\hb^{T}(\nabla\wedge f^{a}) = \sum_{b} e^{b}\wedge(\Omega(e_{b})\cdot e^{a})$, where $f^{a}$ are the background vector fields with $\hb^{T}(f^{a}) = e^{a}$ ($f^{0} = F\,\nabla t$, $f^{1} = \nabla r/F$, $f^{2} = r\,\nabla\vartheta$, $f^{3} = r\sin\vartheta\,\nabla\varphi$). It is linear in the twenty-four components of $\Omega$ and the map is injective (the frame form of the uniqueness of the Levi-Civita connection), so it has exactly one solution. With $\hb^{T}(\nabla r) = F\,e^{1}$, $\hb^{T}(\nabla\vartheta) = e^{2}/r$, $\hb^{T}(\nabla\varphi) = e^{3}/(r\sin\vartheta)$ and $\nabla F = (f'/2F)\nabla r$, the solution is
\[
\Omega(e_{0}) = \frac{f'}{2\sqrt f}\,e_{1}e_{0}, \qquad \Omega(e_{1}) = 0, \qquad \Omega(e_{2}) = \frac{\sqrt f}{r}\,e_{1}\wedge e_{2}, \qquad \Omega(e_{3}) = \frac{\sqrt f}{r}\,e_{1}\wedge e_{3} + \frac{\cot\vartheta}{r}\,e_{2}\wedge e_{3}.
\]
The check solves the torsion equation exactly from the coframe alone at four rational points ($\sqrt f$, $\sin\vartheta$, $\cos\vartheta$ rational, $M = \tfrac12, \tfrac12, \tfrac32, \tfrac13$), finds the linear map of rank $24$, and finds these four bivectors. The covariant derivatives of the frame are then $(e_{c}\cdot D)e_{b} = \Omega(e_{c})\cdot e_{b}$, for instance $(e_{0}\cdot D)e_{0} = (f'/2\sqrt f)\,e_{1}$, $(e_{0}\cdot D)e_{1} = (f'/2\sqrt f)\,e_{0}$, $(e_{2}\cdot D)e_{1} = (\sqrt f/r)\,e_{2}$, $(e_{3}\cdot D)e_{2} = (\cot\vartheta/r)\,e_{3}$: this is the statement $(e_{c}\cdot D)e_{b} = \omega^{a}{}_{b}(e_{c})\,e_{a}$, with $\omega^{a}{}_{b}(e_{c}) = e^{a}\cdot(\Omega(e_{c})\cdot e_{b})$ the components of the bivectors (the check compares all sixteen pairs $(b, c)$ with the table of the Fact).

\emph{The classes.} The pairing class of $\ell = e_{0} + e_{1}$ is $\pv(B) = \ip{B\cdot\ell}{m}$ with $m = e_{2} + I_{5}e_{3}$ in the odd extension $\Cl(2,3)$, $I_{5} = I_{4}e$ central with $I_{5}^{2} = -1$ (\claimref{std-algebra.complexification-is-the-odd-extension}); since $I_{5}$ is central, the scalar product extended bilinearly over $\mathrm{span}\{1, I_{5}\}$ is the symmetric product $\tfrac12(vm + mv)$. From the four bivectors: $(e_{1}e_{0})\cdot\ell = \ell$, so $\Omega(e_{0})\cdot\ell = (f'/2\sqrt f)\,\ell$ and its class is $\ip{\ell}{m} = 0$; $\Omega(e_{1}) = 0$; $\Omega(e_{2})\cdot\ell = (\sqrt f/r)\,e_{2}$, with class $(\sqrt f/r)\ip{e_{2}}{m} = -\sqrt f/r$; and $\Omega(e_{3})\cdot\ell = (\sqrt f/r)\,e_{3}$ (the $e_{2}\wedge e_{3}$ term annihilates $\ell$), with class $(\sqrt f/r)\,I_{5}\,Q(e_{3}) = -I_{5}\sqrt f/r$. By linearity over $\mathrm{span}\{1, I_{5}\}$, with $\Omega(m) = \Omega(e_{2}) + I_{5}\Omega(e_{3})$ and $\Omega(\mb) = \Omega(e_{2}) - I_{5}\Omega(e_{3})$, the classes on the null tetrad are $\pv(\Omega(\ell)) = \pv(\Omega(n)) = \pv(\Omega(m)) = 0$ and $\pv(\Omega(\mb)) = -\tfrac{\sqrt f}{r}(1 - I_{5}^{2}) = -2\sqrt f/r$. The transverse expansion is the trace of $A \mapsto (A\cdot D)\ell$ on the screen, $-\sum_{c = 2,3} e_{c}\cdot(\Omega(e_{c})\cdot\ell) = 2\sqrt f/r = -\pv(\Omega(\mb))$. Along $\ell$ itself, $(\ell\cdot D)\ell = (\Omega(e_{0}) + \Omega(e_{1}))\cdot\ell = (f'/2\sqrt f)\,\ell$, a rescaling. $\Omega(e_{0})$ is a multiple of $e_{01}$ and $\Omega(e_{1}) = 0$, so the transports along the $t$-line and the $r$-line lie in the stabiliser of $\ell$. The check computes every class in $\Cl(2,3)$ with $I_{5}$ from the solved gauge field, at the same four points.

\emph{Horizon limit and the sample value.} At $M = \tfrac12$, $f = 1 - 1/r$ gives $r = 1/(1 - F^{2})$ and $\sqrt f/r = F(1 - F^{2})$, so the transverse class is $-F(1 - F^{2})$, which tends to $0$ as $F \to 0$; the classes of $\Omega(e_{0})$ and $\Omega(e_{1})$ are zero and the others are multiples of $\sqrt f/r$, so all tend to $0$ there. The limit is taken from outside: the static coframe has $e^{1} = dr/F$, so the static gauge is not defined on the horizon locus $F = 0$. At $F = \tfrac35$, $r = \tfrac{25}{16}$ and the class is $-\tfrac35\cdot\tfrac{16}{25} = -\tfrac{48}{125}$.

\emph{In components} (a labelled reading). The components $\omega^{a}{}_{b}(e_{c})$ are the connection forms of Fact \claimref{pairing.imported-schwarzschild-connection-curvature} in the static coframe, the torsion equation is the first structure equation $d\theta^{a} + \omega^{a}{}_{b}\wedge\theta^{b} = 0$, and, read through Fact \claimref{pairing.imported-spin-coefficients-null-tetrad}, the vanishing classes on $\ell$, $n$ and $m$ are the vanishing spin coefficients $\kappa$, $\tau$ and $\sigma$: the congruence is geodesic, $\tau = 0$ and shear-free, with transverse expansion $2\sqrt f/r$. The earlier component checks (the connection forms against the first structure equations and the frame matrices) are kept beside the GTG ones.
